\section{2. Related Work}

\subsection{2.1 Sparse Attention}

\textbf{Longformer.} Beltagy et al. (2020) combine a sliding window of radius 
$W$ with a small set of global tokens. Complexity: $\mathcal{O}(L W)$. 
ClusterFormer captures more global context at similar cost.

\textbf{BigBird.} Zaheer et al. (2020) add random attention to sliding-window. 
Theoretically Turing-complete; empirically similar to Longformer.

\textbf{Sparse Transformer.} Child et al. (2019) use strided and fixed patterns. 
Orthogonal to clustering.

\subsection{2.2 Kernel and Linear Attention}

\textbf{Linear Transformer.} Katharopoulos et al. (2020) replace softmax with 
$\phi(x) = \text{elu}(x) + 1$, achieving $\mathcal{O}(L d^2)$ via 
associative reordering.

\textbf{Performer.} Choromanski et al. (2021) approximate softmax via positive 
orthogonal random features (PORF). Complexity: $\mathcal{O}(L m d)$ with 
$m$ features.

\textbf{Nyströmformer.} Xiong et al. (2021) approximate the attention matrix 
via Nyström landmarks, achieving $\mathcal{O}(L m d)$.

\subsection{2.3 Cluster-Based Attention}

\textbf{Routing Transformer.} Roy et al. (2021) cluster tokens into $k$ groups 
and restrict attention within groups. Complexity: $\mathcal{O}(L \cdot L/k)$. 
ClusterFormer's centroid-to-key cross-attention is fundamentally different.

\textbf{Reformer.} Kitaev et al. (2020) use LSH to approximate attention. Its 
guarantee is probabilistic, not exact.

\textbf{Cluster-Former.} Wang et al. (2021) cluster at sentence granularity for 
document tasks. Not applicable to token-level attention.

\textbf{Pixelfly / Cluster Attention.} Various works explore clustering in vision 
transformers. ClusterFormer targets language.

\subsection{2.4 MoE and Routing}

\textbf{Switch Transformer.} Fedus et al. (2022) route tokens to experts. Similar 
routing mechanism but different purpose.

\textbf{Routing Attention.} Our framework's \texttt{routing} module clusters and routes 
sequentially. ClusterFormer is vectorized and differentiable.

\subsection{2.5 Position of Our Work}

\begin{table}[H]
\centering
\begin{tabular}{lllll}
\toprule
\textbf{Method} & \textbf{Complexity} & \textbf{Global?} & \textbf{Local?} & \textbf{Differentiable?} \\
\midrule
Softmax & $\mathcal{O}(L^2)$ & Yes & Yes & Yes \\
Longformer & $\mathcal{O}(LW)$ & Partial & Yes & Yes \\
Performer & $\mathcal{O}(Lm)$ & Yes & No & Yes \\
Reformer & $\mathcal{O}(L\log L)$ & Yes & No & Approx. \\
Routing & $\mathcal{O}(L^2/k)$ & Partial & No & No (argmin) \\
**ClusterFormer** & $\mathcal{O}(KL)$ & **Yes** & **Yes** & **Yes** \\
\bottomrule
\end{tabular}
\end{table}

ClusterFormer is unique in combining:
1. \textbf{Exact} centroid-to-key attention (not approximate).
2. \textbf{Global + local} blending.
3. \textbf{Fully differentiable} clustering (softmax assignment).
4. \textbf{Multi-architecture} support including GQA and RoPE.

\subsection{2.6 Comparison Table}
